% =============================================================================
%  Pełny opis zatwierdzonych tematów pracy magisterskiej (D2 i P4)
%  Kompilacja: pdflatex magisterkav4.tex  (uruchomić dwukrotnie, dla spisu treści)
%  Wersja: v4 -- po akceptacji tematów przez promotora
% =============================================================================

\documentclass[12pt,a4paper]{article}

\usepackage{polski}
\usepackage[utf8]{inputenc}
\usepackage[T1]{fontenc}

\usepackage{helvet}
\renewcommand{\familydefault}{\sfdefault}

\usepackage[a4paper,left=2.2cm,right=2.2cm,top=2cm,bottom=2cm]{geometry}
\usepackage{setspace}
\setstretch{1.15}

\usepackage{graphicx}
\usepackage{xcolor}
\usepackage{enumitem}
\usepackage{titlesec}
\usepackage{array}
\usepackage{tabularx}
\usepackage{booktabs}
\usepackage{colortbl}
\usepackage{longtable}
\usepackage{tcolorbox}
\tcbuselibrary{skins,breakable}
\usepackage[hidelinks,colorlinks=true,linkcolor=primary,urlcolor=primary]{hyperref}
\usepackage{microtype}

\definecolor{primary}{HTML}{1F3A5F}
\definecolor{accent}{HTML}{C0392B}
\definecolor{lightbg}{HTML}{F4F6FA}
\definecolor{highlightbg}{HTML}{FFF8E1}
\definecolor{highlightborder}{HTML}{F39C12}
\definecolor{border}{HTML}{CCCCCC}
\definecolor{greytxt}{HTML}{555555}
\definecolor{okgreen}{HTML}{1E7B34}

\titleformat{\section}
    {\color{primary}\Large\bfseries}{\thesection}{0.6em}{}
\titleformat{\subsection}
    {\color{primary}\large\bfseries}{\thesubsection}{0.5em}{}
\titleformat{\subsubsection}
    {\color{primary}\normalsize\bfseries}{\thesubsubsection}{0.4em}{}
\titlespacing*{\section}{0pt}{1.4em}{0.7em}
\titlespacing*{\subsection}{0pt}{1.2em}{0.5em}
\titlespacing*{\subsubsection}{0pt}{0.9em}{0.3em}

\setlength{\parindent}{0pt}
\setlength{\parskip}{0.6em}

\setlist[itemize]{leftmargin=1.4em,topsep=0.2em,itemsep=0.2em}
\setlist[enumerate]{leftmargin=1.6em,topsep=0.2em,itemsep=0.2em}

% Nagłówek tematu
\newcommand{\topicheader}[3]{%
    \begin{tcolorbox}[
        colback=primary,colframe=primary,
        boxrule=0pt,arc=0pt,
        left=8pt,right=8pt,top=4pt,bottom=4pt,
        before skip=0.6em, after skip=0.4em
    ]
    \color{white}\bfseries #1 \hfill \normalfont\color{white}Źródło: #2 \quad \bfseries\textbullet\
    \end{tcolorbox}
    \subsection*{#3}
    \addcontentsline{toc}{subsection}{#3}
}

\newenvironment{stacktab}{%
    \arrayrulecolor{border}%
    \begin{tabular}{|>{\columncolor{lightbg}\bfseries\color{primary}\raggedright\arraybackslash}p{4.4cm}|>{\raggedright\arraybackslash}p{11.6cm}|}
    \hline
}{%
    \end{tabular}%
    \arrayrulecolor{black}%
}
\newcommand{\stackrow}[2]{#1 & #2 \\ \hline}

\newcommand{\topicsection}[1]{\subsubsection*{#1}}

\newenvironment{literatura}{%
    \topicsection{Literatura}
    \begin{itemize}[leftmargin=1.6em,labelsep=0.5em]
}{%
    \end{itemize}
}

\newtcolorbox{infobox}[1]{
    colback=highlightbg,
    colframe=highlightborder,
    boxrule=1pt,arc=2pt,
    left=10pt,right=10pt,top=8pt,bottom=8pt,
    title={\color{primary}\bfseries #1},
    coltitle=primary,
    fonttitle=\bfseries,
    breakable
}

\newtcolorbox{okbox}[1]{
    colback=lightbg,
    colframe=okgreen,
    boxrule=1pt,arc=2pt,
    left=10pt,right=10pt,top=8pt,bottom=8pt,
    title={\color{okgreen}\bfseries #1},
    coltitle=okgreen,
    fonttitle=\bfseries,
    breakable
}

% =============================================================================
\begin{document}

% =============================================================================
% STRONA TYTUŁOWA
% =============================================================================
\begin{titlepage}
\null\vspace*{3cm}
\begin{center}
{\color{primary}\Huge\bfseries Szczegółowy opis zatwierdzonych\\[0.3em] tematów pracy magisterskiej}\\[1.4em]
{\color{greytxt}\Large Temat D2 -- Dawid Olko \quad$\bullet$\quad Temat P4 -- Piotr Smoła}\\[0.8em]
{\color{greytxt}\large Rozszerzenie pracy inżynierskiej:\\ System rekomendacji produktów oparty na metodach hybrydowych}
\end{center}

\vspace{2cm}

Niniejszy dokument stanowi rozwinięcie wcześniejszej listy propozycji tematów. Tematy to: \textbf{Pre-trenowane modele językowe (LLM) otwarto-źródłowe a klasyczne metody rekomendacji} (autor: Dawid Olko) oraz \textbf{Wnioskowanie przedziałowe (logika rozmyta typu drugiego, IT2-FLS) z biblioteką \texttt{pyit2fls}} (autor: Piotr Smoła).

\vspace{0.4em}

W dokumencie przedstawiono: (1) analizę obecnego stanu aplikacji \emph{SmartRecommender} -- co już działa i co stanowi punkt wyjścia; (2) pełny, uszczegółowiony opis obu tematów -- co, gdzie i w jaki sposób będzie realizowane; (3) uwzględnienie wskazówek (język pracy i promptów, polskie modele jako punkt odniesienia, włączenie ANFIS, fundament wspólny obu prac); (4) przegląd rekomendowanych modeli LLM wraz z uzasadnieniem doboru; (5) analizę literatury dotyczącej minimalnej liczby użytkowników wymaganej dla wiarygodności statystycznej testów.

\vfill
\begin{center}
{\color{greytxt}Uniwersytet Rzeszowski \quad$\bullet$\quad rok akademicki 2025/2026}
\end{center}
\end{titlepage}

\tableofcontents
\newpage

% =============================================================================
% 1. USTALENIA Z PROMOTOREM
% =============================================================================
\section{Ustalenia z promotorem}

Poniżej ustalono dodatkowo cztery kwestie, które zostały uwzględnione w niniejszym opisie.

\begin{enumerate}
\item \textbf{Język pracy i promptów.} Praca magisterska (temat D2) zostanie napisana w języku polskim. Prompty kierowane do modeli językowych zostaną przygotowane w języku angielskim -- jest to język, w którym modele uniwersalne osiągają najwyższą jakość, a jednocześnie utrzymuje to spójność z literaturą. Główną oś porównawczą stanowią modele anglojęzyczne (Llama, Mistral, Qwen i inne), natomiast polskie modele otwarto-źródłowe (Bielik, PLLuM) zostają włączone jako ciekawy punkt odniesienia oraz przykład lokalnych inicjatyw badawczych.

\item \textbf{Polskie systemy i modele.} Zgodnie ze wskazówką -- ,,skoro istnieją polskie systemy, to należy je uwzględnić'' -- temat D2 obejmuje przegląd i empiryczną ocenę polskich modeli Bielik oraz PLLuM. Stanowi to wyraźny lokalny wkład pracy.

\item \textbf{Rozszerzenie tematu P4 o ANFIS.} Zasugerowana możliwość włączenia ANFIS (,,można do badania włączyć ANFIS i porównać te trzy systemy''). Trzon tematu P4 stanowi porównanie systemu typu pierwszego (obecny Mamdani-TSK) z systemem typu drugiego przedziałowym (IT2-FLS, biblioteka \texttt{pyit2fls}). Po analizie wykonalności (sekcja ,,Analiza wykonalności ANFIS'') ANFIS został potraktowany jako \emph{rozszerzenie warunkowe} -- ze względu na problem etykiet treningowych, którego para T1 kontra IT2-FLS nie ma. Ostateczny status ANFIS musi zostać uzgodniony i zatwierdzony bądź odrzucony.

\item \textbf{Wiarygodność statystyczna i liczba użytkowników.} Liczba użytkowników wymagana dla istotności statystycznej zależy od minimalnych wymagań stosowanych testów -- zwykle wystarcza kilkanaście do kilkudziesięciu użytkowników -- sprawdzenie liczb spotykanych w literaturze systemów rekomendacyjnych. Przegląd ten zawiera rozdział~\ref{sec:liczebnosc}.
\end{enumerate}

% =============================================================================
% 2. STAN OBECNY APLIKACJI
% =============================================================================
\section{Stan obecny aplikacji \emph{SmartRecommender}}

Praca inżynierska zaowocowała w pełni działającym systemem e-commerce zintegrowanym w architekturze \textbf{Django + Django REST Framework + React + PostgreSQL}, z modułem rekomendacji obejmującym dziewięć metod. Praca magisterska nie buduje sklepu od nowa -- rozszerza istniejący system o warstwę badawczo-eksperymentalną. Poniższy opis odzwierciedla faktyczny stan repozytorium i posłuży jako punkt odniesienia dla obu tematów.

\subsection*{Architektura i model danych}

Backend (katalog \texttt{backend/home/}) opiera się na rozszerzonym modelu użytkownika z kontrolą dostępu opartą na rolach (admin / klient) oraz na zestawie modeli domenowych. Kluczowe encje to: \texttt{Product} (nazwa, cena, opis, kategorie, tagi -- opis tekstowy jest wykorzystywany w filtracji opartej na zawartości), \texttt{Category}, \texttt{Order} i \texttt{OrderProduct} (oś czasowa zamówień), \texttt{Opinion} (oceny 1--5 oraz treść tekstowa, z ograniczeniem unikalności para użytkownik--produkt), \texttt{Specification} (parametry techniczne) oraz \texttt{Complaint} (reklamacje).

Osobną grupę stanowią modele warstwy metod, których nazwy tabel rozpoczynają się od przedrostka \texttt{method\_}: \texttt{UserInteraction} (zdarzenia: \emph{view}, \emph{click}, \emph{add\_to\_cart}, \emph{purchase}, \emph{favorite}), \texttt{ProductSimilarity} (zapisane podobieństwa produkt--produkt, z rozróżnieniem typu: \emph{collaborative} lub \emph{content\_based}), \texttt{ProductAssociation} (reguły asocjacyjne z miarami support, confidence, lift), \texttt{SentimentAnalysis} i \texttt{ProductSentimentSummary} (wynik analizy sentymentu opinii), \texttt{PurchaseProbability}, \texttt{UserPurchasePattern}, \texttt{SalesForecast}, \texttt{ProductDemandForecast}, \texttt{RiskAssessment} oraz \texttt{UserProductRecommendation}. Taki podział oznacza, że wyniki metod są \emph{materializowane} w bazie -- nie są przeliczane przy każdym żądaniu -- co jest istotne dla powtarzalności eksperymentów.

\subsection*{Zaimplementowane metody rekomendacji}

\begin{tabularx}{\linewidth}{|c|X|l|}
\hline
\rowcolor{primary}
\textcolor{white}{\textbf{Lp.}} &
\textcolor{white}{\textbf{Metoda i implementacja}} &
\textcolor{white}{\textbf{Dane}} \\ \hline
1 & Filtracja oparta na zawartości -- TF-IDF, ważone cechy (kategorie, tagi, słowa kluczowe z opisu, kategoria cenowa), podobieństwo kosinusowe; klasa \texttt{CustomContentBasedFilter} & jawne, atrybuty \\ \hline
2 & Filtracja kolaboratywna item-based -- Adjusted Cosine Similarity (Sarwar i in. 2001), wyniki w \texttt{ProductSimilarity} & ukryte, transakcje \\ \hline
3 & Wnioskowanie rozmyte typu Mamdaniego -- klasa \texttt{SimpleFuzzyInference}, 6 reguł, zmienne: cena / jakość / popularność & mieszane \\ \hline
4 & Wyszukiwanie rozmyte -- klasa \texttt{CustomFuzzySearch}: odległość edycyjna, podobieństwo trigramowe i znakowe & tekst zapytań \\ \hline
5 & Analiza sentymentu opinii -- klasa \texttt{CustomSentimentAnalysis}: metoda leksykonowa z negacją, intensyfikatorami i analizą bigramów & tekst opinii \\ \hline
6 & Łańcuchy Markowa -- klasa \texttt{CustomMarkovChain}, przejścia między produktami / kategoriami & sekwencyjne \\ \hline
7 & Klasyfikator naiwny Bayesa -- klasa \texttt{CustomNaiveBayes}, predykcja prawdopodobieństwa zakupu & mieszane \\ \hline
8 & Reguły asocjacyjne -- klasa \texttt{CustomAssociationRules}: Apriori z reprezentacją bitmapową, parametry \texttt{min\_support}, \texttt{min\_confidence} & transakcje \\ \hline
9 & Analiza sezonowości -- komponent temporalny prognoz popytu i sprzedaży & czasowe \\ \hline
\end{tabularx}

Wszystkie metody zaimplementowano autorsko w pliku \texttt{custom\_recommendation\_engine.py} (silnik rozmyty -- w osobnym \texttt{fuzzy\_logic\_engine.py}), bez korzystania z gotowych bibliotek rekomendacyjnych. Każda metoda ma własny zestaw widoków API (np. \texttt{recommendation\_views.py}, \texttt{association\_views.py}, \texttt{sentiment\_views.py}) oraz panel diagnostyczny w interfejsie administratora, co umożliwia niezależne uruchomienie i podgląd działania -- jest to mocna podstawa pod eksperymenty porównawcze.

\subsection*{Zasilanie danymi -- obecny seeder}

Aplikacja zawiera komendę \texttt{python manage.py seed} (plik \texttt{seed.py}), zasilającą bazę danymi demonstracyjnymi sklepu komputerowego: kategorie, użytkownicy, produkty z realnymi opisami i specyfikacjami, opinie, zamówienia, reklamacje, a na końcu przeliczenie tabel pochodnych (sentyment, reguły asocjacyjne, podobieństwa, prognozy). Dane są jednak częściowo losowe i w niewielkiej liczbie -- wystarczające do demonstracji systemu, lecz niewystarczające do rzetelnych eksperymentów. Stąd potrzeba fundamentu wspólnego opisanego w rozdziale~\ref{sec:fundament}.

% =============================================================================
% 3. FUNDAMENT WSPÓLNY
% =============================================================================
\section{Fundament wspólny obu prac}
\label{sec:fundament}

Poniżej przedstawiono zakres tego fundamentu w formie, którą obie prace współdzielą. Jest to pierwszy etap realizowany wspólnie przez obu autorów; dopiero na nim budowane są osobne rozdziały eksperymentalne.

\begin{enumerate}
\item \textbf{Badawczy seeder danych} -- nowa komenda \texttt{python manage.py seed\_research}, działająca obok istniejącego \texttt{seed} (bez kasowania go, bez zmian schematu bazy). Seeder generuje parametryzowaną liczbę użytkowników z \emph{personami zakupowymi} (np. gracz, biuro, twórca treści, kupujący prezent, entuzjasta), spójne tematycznie koszyki, oś czasową zamówień oraz strumienie danych jawnych (opinie) i ukrytych (interakcje). Stały \emph{seed} liczbowy gwarantuje powtarzalność wyników. Parametry siły sygnału i poziomu szumu są jawne i kontrolowane.

\item \textbf{Polecenie \texttt{prepare\_research\_db}} -- po zasileniu bazy zawsze wykonuje to samo przeliczenie tabel pochodnych (macierz podobieństw, reguły asocjacyjne, agregaty sentymentu), tak aby każdy eksperyment startował z identycznego stanu bazy.

\item \textbf{Wspólny moduł metryk} (np. \texttt{home/experiments/metrics.py}) -- ten sam podział temporalny (ostatni ustalony procent osi czasu = zbiór testowy) oraz te same metryki rankingowe: Precision@K, Recall@K, nDCG@K, dla raz ustalonego $K$.

\item \textbf{Miejsce zapisu wyników} -- model \texttt{ExperimentRun} (parametry konfiguracji, metryki, czas wykonania, znacznik czasu) lub eksport do plików CSV -- aby wyniki nie ginęły w logach konsoli i były gotowe do analizy statystycznej.
\end{enumerate}

Dzięki wspólnemu fundamentowi temat D2 i temat P4 korzystają z tego samego zbioru danych, tego samego podziału i tych samych metryk -- co umożliwia bezpośrednie porównanie wyników w sekcji wspólnej obu prac.

\newpage

% =============================================================================
% 4. TEMAT D2
% =============================================================================
\section{Temat D2 -- pełny opis}

\topicheader{Temat D2}{autorski (Dawid Olko)}{Pre-trenowane modele językowe (LLM) otwarto-źródłowe a klasyczne metody rekomendacji}

\topicsection{Cel i pytanie badawcze}
Temat odpowiada na aktualne pytanie literatury rekomendacyjnej: czy nowoczesne, otwarto-źródłowe duże modele językowe (LLM) mogą realnie konkurować z klasycznymi metodami rekomendacji, a jeśli tak -- w jakim trybie wykorzystania, jakim kosztem i z jakim ryzykiem błędów. Praca porównuje rekomendacje generowane przez dostępne klasyczne metody z pracy inżynierskiej z rekomendacjami uzyskanymi z modeli LLM, w jednolitym protokole ewaluacji.

\topicsection{Trzy paradygmaty wykorzystania LLM}
Praca eksploruje trzy sposoby włączenia modelu językowego do systemu rekomendacji -- od najprostszego do najbardziej złożonego:

\begin{enumerate}
\item \textbf{Bezpośrednia rekomendacja / re-ranking (prompting).} Model otrzymuje w prompcie historię zakupów użytkownika oraz \emph{listę produktów-kandydatów} -- zawężoną wcześniej przez tanią metodę klasyczną (filtracja oparta na treści lub kolaboratywna), aby model nie musiał przeszukiwać całego katalogu. Zadaniem modelu jest uporządkowanie listy. Badane są warianty zero-shot i few-shot. Zgodnie z uwagą, prompty są przygotowane w języku angielskim.

\item \textbf{Reprezentacje wektorowe opisów (embeddings).} Model językowy (lub dedykowany model embeddujący) oblicza wektorowe reprezentacje opisów produktów i treści opinii. Reprezentacje są indeksowane w lekkiej bazie wektorowej, a rekomendacja sprowadza się do wyszukiwania najbliższych sąsiadów. Ten paradygmat naturalnie radzi sobie z problemem zimnego startu produktu -- nowy produkt z opisem jest rekomendowalny natychmiast, mimo braku historii zakupów.

\item \textbf{Retrieval-Augmented Generation (RAG).} Kontekst przekazywany do modelu generującego ranking jest wzbogacany o powiązane opinie i specyfikacje techniczne, pobierane z bazy wektorowej. Hipoteza: RAG ogranicza halucynacje względem prostego promptingu, ale podnosi opóźnienie odpowiedzi.
\end{enumerate}

\begin{infobox}{Spójność tematów D2 i P4}
Oba tematy badają tę samą rzecz -- czy nowocześniejsze, bardziej złożone podejście realnie poprawia jakość rekomendacji, czy jedynie podnosi koszt. D2 zadaje to pytanie dla modeli językowych, P4 -- dla formalnego modelowania niepewności (logika rozmyta typu drugiego, ANFIS). Oba korzystają z tego samego seederа badawczego, tego samego podziału temporalnego i tych samych metryk, co pozwala zestawić wyniki w jednym wspólnym rozdziale.
\end{infobox}

\topicsection{Zakres prac implementacyjnych}
Powstaje moduł rekomendacji oparty na LLM, działający równolegle do istniejących modułów -- bez ingerencji w działający system. Moduł składa się z trzech klas odpowiadających trzem paradygmatom:

\begin{itemize}
\item \textbf{Klasa promptująca} -- buduje prompt zawierający historię zakupów i listę kandydatów, wywołuje model, parsuje odpowiedź. Odpowiedź modelu jest walidowana strukturalnie (np. Pydantic v2); w razie niepoprawnej odpowiedzi system stosuje fallback do rankingu bazowego.
\item \textbf{Klasa embeddująca} -- liczy reprezentacje wektorowe opisów i opinii, indeksuje je, obsługuje zapytania semantyczne (wyszukiwanie kNN).
\item \textbf{Klasa RAG} -- pipeline pobierający kontekst i wzbogacający prompt przed wywołaniem modelu generującego ranking.
\end{itemize}

Wszystkie klasy korzystają wyłącznie z modeli otwarto-źródłowych, co zapewnia pełną reprodukowalność i niezależność od polityk komercyjnych dostawców.

\topicsection{Stack technologiczny}
\begin{stacktab}
\stackrow{Modele LLM uniwersalne}{Llama 4 / 4 Scout, Mistral (Medium 3.5 / Large 3, Apache 2.0), Qwen 3.5, GLM-5.1, DeepSeek V4 -- patrz rozdział~\ref{sec:modele}}
\stackrow{Modele LLM lokalne}{Mistral Small (24B), Gemma 4, Phi-4 (14B), Qwen 3.5 w mniejszych wariantach -- działają na jednej karcie GPU}
\stackrow{Modele LLM polskojęzyczne}{Bielik v3 (7B i 11B, SpeakLeash, Apache 2.0), PLLuM (8B--70B, konsorcjum OPI PIB) -- jako punkt odniesienia}
\stackrow{Reprezentacje wektorowe}{sentence-transformers; modele wielojęzyczne BGE / E5 (porównanie modelu wielojęzycznego z anglojęzycznym)}
\stackrow{Baza wektorowa}{pgvector (rozszerzenie PostgreSQL -- bez nowych usług w stacku) lub Qdrant / ChromaDB w trybie lokalnym}
\stackrow{Framework RAG}{LangChain lub LlamaIndex}
\stackrow{Inferencja}{Ollama, llama.cpp, vLLM lub Hugging Face Transformers}
\stackrow{Backend}{Django + DRF -- nowy moduł rekomendacji LLM równolegle do istniejących}
\end{stacktab}

\topicsection{Innowacyjność}
Wkład pracy obejmuje cztery elementy. Po pierwsze, systematyczne porównanie modeli LLM różnej skali i pochodzenia z klasycznymi metodami rekomendacji na tym samym zbiorze danych e-commerce. Po drugie, analiza kompromisu między jakością rekomendacji a kosztem obliczeniowym (opóźnienie odpowiedzi, zapotrzebowanie na pamięć GPU). Po trzecie, ocena współczynnika halucynacji -- częstości, z jaką model rekomenduje produkty nieistniejące w katalogu lub przypisuje im fałszywe atrybuty; jest to metryka rzadko raportowana, a krytyczna dla wdrożeń produkcyjnych. Po czwarte, badanie wpływu języka oraz ocena polskich modeli Bielik i PLLuM jako lokalnie dostępnej, taniej alternatywy -- co stanowi wyraźny lokalny wkład pracy.

\topicsection{Hipotezy badawcze}
\begin{itemize}
\item \textbf{H1.} Model LLM operujący na liście kandydatów dostarczonej przez tanią metodę zawężającą osiąga wyższą metrykę nDCG@K niż ta metoda samodzielnie.
\item \textbf{H2.} Jakość rekomendacji rośnie wraz ze skalą modelu, ale punkt nasycenia leży poniżej największych dostępnych modeli -- stosowanie modeli największych nie opłaca się, jeśli mniejsze dają porównywalny wynik.
\item \textbf{H3.} Polski model Bielik v3 11B operujący na polskich opisach produktów osiąga jakość niegorszą niż wielokrotnie większe modele uniwersalne, przy znacząco niższym koszcie obliczeniowym.
\item \textbf{H4.} Pipeline RAG zmniejsza współczynnik halucynacji względem prostego promptingu, lecz podnosi opóźnienie odpowiedzi.
\end{itemize}

\topicsection{Plan eksperymentu}
Przygotowanie reprezentatywnego zbioru zapytań rekomendacyjnych z badawczego zbioru danych. Generacja rekomendacji przez wybrane modele otwarto-źródłowe (od czterech do sześciu modeli reprezentujących różne klasy wielkości i pochodzenia, w tym co najmniej jeden polski). Generacja rekomendacji przez wszystkie metody klasyczne oraz wybrane warianty hybrydowe (LLM nad metodą klasyczną). Pomiar metryk Precision@K, nDCG@K, współczynnika halucynacji, średniego opóźnienia odpowiedzi i kosztu zasobowego. Analiza statystyczna -- test Friedmana dla porównania wielu metod, z testem post-hoc (np. Nemenyiego) dla par metod. Jako uzupełnienie jakościowe -- ocena ekspercka przez autorów na niewielkiej, kontrolnej próbie przypadków (bez badania z udziałem użytkowników -- szczegóły w rozdziale~\ref{sec:liczebnosc}).

\topicsection{Koszty modeli i nakład pracy na wdrożenie}
Świadomie wybieram wyłącznie modele otwarto-źródłowe, więc \textbf{za same modele nie płacę nic (oprócz prądu elektrycznego)} -- pobieram je bezpłatnie z platformy Hugging Face i uruchamiam lokalnie na własnym sprzęcie. W odróżnieniu od komercyjnych API (OpenAI, Anthropic, Google), gdzie płaci się za każde tysiąc tokenów, tutaj nie ma żadnych opłat za zapytanie ani limitów -- mogę wykonać dowolnie wiele przebiegów eksperymentu, co jest istotne dla powtarzalności badania. Jedyny realny koszt to prąd i czas pracy mojego komputera.

Pozostają trzy potencjalne, opcjonalne źródła kosztów -- żadne nie jest konieczne do realizacji pracy:

\begin{itemize}
\item \textbf{Modele zbyt duże na moją kartę graficzną} (rzędu 70B parametrów i większe). Jeśli zdecyduję się je włączyć do porównania, mogę skorzystać z usług hostingu inferencji rozliczanych za użycie (np. Together AI, Groq, OpenRouter) -- koszt jednego pełnego przebiegu eksperymentu na takim modelu to zwykle kilka--kilkanaście złotych, a całość mieści się w kilkudziesięciu złotych. Alternatywnie -- bezpłatnie -- uczelniany klaster obliczeniowy lub darmowy limit GPU w środowisku Google Colab.
\item \textbf{Wynajem GPU w chmurze} (np. RunPod, Vast.ai) -- kilka złotych za godzinę, gdyby mój sprzęt okazał się niewystarczający. Przy planowanym zestawie 4--6 modeli mieszczących się lokalnie nie przewiduję takiej potrzeby.
\item \textbf{Modele polskie} -- Bielik (licencja Apache 2.0) i PLLuM są w pełni darmowe i uruchamiają się lokalnie.
\end{itemize}

\textbf{Czas nauki.} Nie jest wymagane trenowanie ani dostrajanie modeli (fine-tuning) -- modele wykorzystuję w trybie gotowej inferencji (zero-shot i few-shot), co eliminuje najtrudniejszą i najbardziej czasochłonną część pracy z modelami językowymi. Realny nakład na opanowanie potrzebnych narzędzi szacuję następująco: uruchomienie modelu lokalnie przez Ollamę -- kwestia kilku godzin (instalacja, pobranie modelu, pierwsze zapytanie); konstrukcja promptów i walidacja odpowiedzi -- kilka dni iteracyjnej pracy; podstawy RAG i bazy wektorowej (pgvector, biblioteka LangChain lub LlamaIndex) -- około tygodnia. Najwięcej czasu pochłonie nie nauka narzędzi, lecz właściwe eksperymenty: dobór promptów, przebiegi pomiarowe i analiza wyników. Całość mieści się w typowym harmonogramie pracy magisterskiej i nie wymaga wcześniejszego doświadczenia z modelami językowymi.

\begin{infobox}{W skrócie -- koszty}
Modele: \textbf{0 zł} (otwarto-źródłowe, uruchamiane lokalnie). Brak opłat za zapytania, brak limitów. Opcjonalny koszt tylko gdybym chciał włączyć modele największej skali -- rzędu kilkudziesięciu złotych za hosting rozliczany za użycie, albo bezpłatnie przez klaster uczelni lub Google Colab. Trening modeli nie jest potrzebny.
\end{infobox}

\topicsection{Metodyka badania -- jak dane są porównywane i jak powstają wnioski}
Praca ma charakter badawczy -- jej trzon stanowi systematyczne, mierzalne porównanie metod i modeli, zakończone wnioskami popartymi analizą statystyczną. Poniżej opisano, jak to porównanie jest przeprowadzane krok po kroku.

\textbf{Krok 1 -- punkt odniesienia (ground truth).} Aby w ogóle dało się ocenić, czy rekomendacja jest trafna, potrzebny jest obiektywny wzorzec ,,prawidłowej odpowiedzi''. Stosuję \emph{podział temporalny}: oś czasu zakupów każdego użytkownika dzielona jest na część uczącą (wcześniejsze ok. 80\% historii) i część testową (ostatnie ok. 20\%). Metoda lub model ,,widzi'' wyłącznie część uczącą i na jej podstawie generuje ranking; ranking ten jest następnie konfrontowany z produktami, które użytkownik \emph{faktycznie kupił} w odciętej części testowej. Produkt obecny w przyszłych zakupach to trafienie. Podział jest deterministyczny (stały \emph{seed}), więc każde porównanie jest powtarzalne i nie wymaga udziału ludzi.

\textbf{Krok 2 -- wspólny protokół.} Wszystkie metody klasyczne i wszystkie modele LLM oceniane są w \emph{identycznych} warunkach: ten sam badawczy zbiór danych, ten sam podział temporalny, ta sama wartość $K$, te same metryki. To warunek rzetelności -- bez niego porównanie nie miałoby wartości naukowej. Protokół realizuje wspólny moduł metryk (\texttt{home/experiments/metrics.py}) z fundamentu wspólnego.

\textbf{Krok 3 -- metryki.} Dla każdego użytkownika i każdej metody (lub modelu) obliczane są:
\begin{itemize}
\item \textbf{Precision@K} -- jaka część z $K$ poleconych produktów okazała się trafiona;
\item \textbf{Recall@K} -- jaką część rzeczywistych przyszłych zakupów udało się odzyskać;
\item \textbf{nDCG@K} -- czy trafione produkty znalazły się wysoko w rankingu (pozycja ma znaczenie -- trafienie na miejscu 1. liczy się więcej niż na 10.);
\item \textbf{współczynnik halucynacji} -- odsetek rekomendacji wskazujących produkty spoza katalogu lub z fałszywymi atrybutami (metryka specyficzna dla LLM);
\item \textbf{opóźnienie odpowiedzi i koszt zasobowy} -- średni czas wygenerowania rekomendacji oraz zapotrzebowanie na pamięć GPU.
\end{itemize}

\textbf{Krok 4 -- trzy osie porównania.} Badanie zestawia wyniki na trzech płaszczyznach:
\begin{itemize}
\item \emph{LLM kontra metody klasyczne} -- np. Bielik kontra filtracja kolaboratywna kontra reguły asocjacyjne, na tych samych użytkownikach (główne pytanie pracy);
\item \emph{LLM kontra LLM} -- porównanie modeli między sobą: Qwen, Mistral, Bielik, Phi-4 i inne (weryfikacja hipotez H2 i H3 o wpływie skali i pochodzenia modelu);
\item \emph{paradygmat kontra paradygmat} -- prompting kontra embeddings kontra RAG (weryfikacja hipotezy H4).
\end{itemize}

\textbf{Krok 5 -- analiza statystyczna.} Sama tabela średnich nie wystarcza -- różnica może być przypadkowa. Każda metoda daje wektor wyników (po jednej wartości metryki na każdego użytkownika testowego). \textbf{Test Friedmana} sprawdza, czy różnice między wszystkimi porównywanymi metodami są istotne statystycznie (przy poziomie istotności $\alpha = 0{,}05$). Jeśli test wykaże istotność, \textbf{test post-hoc Nemenyiego} wskazuje, które konkretnie pary metod różnią się istotnie. Pozwala to przejść od stwierdzenia ,,model A wygląda lepiej'' do ,,model A jest istotnie lepszy od modelu B (p = 0{,}03)''. Dla porównań par metod stosowany jest dodatkowo test Wilcoxona dla par powiązanych.

\textbf{Krok 6 -- wnioski.} Wyniki podsumowuje ranking metod według jakości oraz wykres kompromisu \emph{jakość kontra koszt} (front Pareto: nDCG@K na jednej osi, opóźnienie odpowiedzi na drugiej). Wniosek pracy ma postać konkretnej, popartej liczbami rekomendacji -- na przykład: ,,dla sklepu tej skali model Bielik 11B w trybie promptingu osiąga około 90\% jakości największego modelu uniwersalnego przy ułamku kosztu obliczeniowego, a pipeline RAG nie poprawia wyniku w sposób istotny statystycznie''. Tak sformułowany wniosek odpowiada wprost na pytanie badawcze pracy i jest weryfikowalny.

\begin{infobox}{Aspekt badawczy w skrócie}
Praca nie jest wyłącznie wdrożeniem modeli LLM -- jest \textbf{badaniem porównawczym}. Mierzy te same metryki dla wszystkich metod w identycznych warunkach, sprawdza istotność różnic testami statystycznymi (Friedman, post-hoc Nemenyi, Wilcoxon) i kończy się wnioskami popartymi liczbami: który model i który paradygmat daje najlepszy stosunek jakości do kosztu. To właśnie stanowi naukową wartość tematu D2.
\end{infobox}

\topicsection{Przeprowadzenie testów i panel administratora}
Aby eksperymenty były wygodne do uruchamiania i przejrzyste w odbiorze, warstwa badawcza tematu D2 zostanie zintegrowana z istniejącym panelem administratora -- analogicznie do obecnych paneli diagnostycznych poszczególnych metod. Powstanie nowa zakładka \emph{,,Eksperymenty LLM''}, a w niej:

\begin{itemize}
\item \textbf{Uruchamianie eksperymentu z panelu} -- formularz pozwalający wybrać badane modele, paradygmat (prompting / embeddings / RAG), wartość $K$ oraz podzbiór użytkowników testowych, a następnie uruchomić przebieg jednym przyciskiem. Eksperyment wykonuje się w tle (zadanie asynchroniczne), a panel pokazuje pasek postępu -- nie trzeba uruchamiać niczego z terminala.
\item \textbf{Tabela wyników} -- każdy zakończony przebieg trafia jako wiersz do modelu \texttt{ExperimentRun} i jest wyświetlany w panelu: nazwa modelu, paradygmat, metryki Precision@K, nDCG@K, współczynnik halucynacji, średnie opóźnienie odpowiedzi, znacznik czasu. Wyniki można sortować i filtrować.
\item \textbf{Wykresy porównawcze} -- panel rysuje zestawienie metryk dla wszystkich przebiegów (np. wykres słupkowy nDCG@K dla kolejnych modeli oraz wykres jakość kontra opóźnienie), co daje natychmiastowy obraz, który model wypada najlepiej.
\item \textbf{Eksport do CSV} -- przycisk \emph{,,Eksportuj wyniki''} pobiera wszystkie przebiegi (lub wybrane) w postaci pliku CSV gotowego do dalszej analizy statystycznej w środowisku zewnętrznym (Python, R, arkusz kalkulacyjny). Eksport CSV jest też dostępny bezpośrednio przez endpoint API, na wypadek pracy wsadowej.
\item \textbf{Podgląd pojedynczej rekomendacji} -- dla wybranego użytkownika panel pokazuje obok siebie listę kandydatów, ranking zwrócony przez model oraz -- przy paradygmacie RAG -- fragmenty kontekstu (opinie, specyfikacje), które trafiły do promptu. Ułatwia to diagnostykę i jest atrakcyjnym materiałem na obronę.
\end{itemize}

Takie podejście oznacza, że całe badanie da się przeprowadzić i przeglądać ,,klikając'' w panelu, a wyniki w każdej chwili wyeksportować do CSV -- bez grzebania w konsoli.

\topicsection{Widoczność zmian po stronie frontendu}
Tak -- zmiany będą widoczne w interfejsie, i to w dwóch miejscach. Po pierwsze, w \textbf{panelu administratora} pojawi się opisana wyżej zakładka eksperymentów. Po drugie, na \textbf{stronie sklepu} rekomendacje generowane przez modele LLM zostaną podpięte jako kolejna, równoprawna opcja w istniejącym mechanizmie wyboru algorytmu rekomendacji.

Obecnie aplikacja pozwala wskazać aktywny algorytm rekomendacji (model \texttt{RecommendationSettings} z wariantami: filtracja kolaboratywna, oparta na zawartości, logika rozmyta), a sekcje rekomendowanych produktów na stronach renderują się zgodnie z tym wyborem. W ramach tematu D2 lista dostępnych algorytmów zostanie rozszerzona o opcję opartą na LLM -- po jej wybraniu sekcje rekomendacji (np. ,,Polecane dla Ciebie'', ,,Produkty dopasowane'') będą wyświetlać produkty uszeregowane przez model językowy. Dzięki temu admin może w panelu przełączyć sklep na rekomendacje LLM dokładnie tak samo, jak dziś przełącza między metodami klasycznymi, a efekt jest natychmiast widoczny dla klienta na stronie. Pozwala to też pokazać działanie na żywo podczas obrony -- przełączenie algorytmu i porównanie, jak zmienia się zawartość sekcji rekomendacji.

\topicsection{Sposób podpięcia modeli i miejsce konfiguracji promptów}
Modele językowe podpinam do aplikacji przez \emph{warstwę inferencji} -- lokalny serwer (Ollama lub vLLM) udostępniający model pod adresem HTTP, z którym backend Django komunikuje się tak samo, jak z dowolnym innym serwisem. Adres serwera i nazwa modelu trafiają do konfiguracji aplikacji (plik ustawień / zmienne środowiskowe), więc zmiana modelu nie wymaga modyfikacji kodu -- wystarczy wskazać inny model w konfiguracji lub, docelowo, w panelu administratora.

Prompty \textbf{nie są zaszyte na sztywno w kodzie}. Tekst promptu rozkłada się na dwie części. Część stała -- \emph{szablon promptu} (instrukcja zadania, format oczekiwanej odpowiedzi, ewentualne przykłady few-shot) -- przechowywana jest jako wersjonowany szablon: w pierwszej wersji w pliku konfiguracyjnym repozytorium, docelowo w bazie danych, co umożliwi edycję szablonu wprost z panelu administratora bez ponownego wdrażania aplikacji. Część zmienna -- historia zakupów użytkownika i lista produktów-kandydatów -- jest wstawiana do szablonu dynamicznie przy każdym żądaniu. Takie rozdzielenie ma trzy zalety: pozwala porównywać różne wersje promptów w eksperymencie (a wynik każdego przebiegu w tabeli \texttt{ExperimentRun} zapisuje, której wersji promptu użyto -- co gwarantuje powtarzalność), umożliwia łatwą edycję promptów bez dotykania kodu oraz utrzymuje zgodność z ustaleniem językowym -- szablon redaguję po angielsku, mimo że dane produktów pozostają w języku polskim.

\begin{literatura}
\item Wu, L. i in. (2024). A Survey on Large Language Models for Recommendation. \textit{arXiv:2305.19860}.
\item Hou, Y. i in. (2024). Large Language Models are Zero-Shot Rankers for Recommender Systems. \textit{ECIR '24}.
\item Bao, K. i in. (2023). TALLRec: An Effective and Efficient Tuning Framework to Align LLM with Recommendation. \textit{RecSys '23}.
\item Lewis, P. i in. (2020). Retrieval-Augmented Generation for Knowledge-Intensive NLP Tasks. \textit{NeurIPS}.
\item Reimers, N., Gurevych, I. (2019). Sentence-BERT: Sentence Embeddings using Siamese BERT-Networks. \textit{EMNLP}.
\item Ociepa, K. i in. (2025). Advancing Polish Language Modeling through Tokenizer Optimization in the Bielik v3 Series. \textit{arXiv}.
\end{literatura}

\newpage

% =============================================================================
% 5. REKOMENDOWANE MODELE LLM
% =============================================================================
\section{Rekomendowane modele LLM dla tematu D2}
\label{sec:modele}

Dobór modeli do badania kieruje się czterema kryteriami: (1) \textbf{licencja otwarto-źródłowa} pozwalająca na użycie badawcze i reprodukcję; (2) \textbf{możliwość uruchomienia lokalnego} -- przynajmniej części modeli -- bez zależności od komercyjnego API; (3) \textbf{reprezentatywność skali} -- modele małe, średnie i duże, aby zbadać hipotezę o punkcie nasycenia; (4) \textbf{aktualność} -- modele stanowiące stan wiedzy w pierwszej połowie 2026 roku. Poniżej zestaw rekomendowany, podzielony na trzy osie porównawcze.

\subsection*{Oś główna -- modele uniwersalne anglojęzyczne}

\begin{tabularx}{\linewidth}{|>{\bfseries\color{primary}}p{3.0cm}|p{2.2cm}|X|}
\hline
\rowcolor{primary}
\textcolor{white}{\textbf{Model}} & \textcolor{white}{\textbf{Skala / licencja}} & \textcolor{white}{\textbf{Uzasadnienie doboru}} \\ \hline
Qwen 3.5 & wiele wariantów; Apache 2.0 & Najlepszy stosunek jakości do kosztu i otwartości w 2026 r.; dostępny w wariantach od kilku do kilkuset miliardów parametrów -- pozwala badać hipotezę H2 (punkt nasycenia) w obrębie jednej rodziny modeli. \\ \hline
Llama 4 / 4 Scout & MoE; licencja Meta & Najszerzej wdrażana rodzina modeli otwartych; wariant Scout dobrze radzi sobie z długim kontekstem -- istotne dla RAG i długich list kandydatów. \\ \hline
Mistral (Medium 3.5 / Large 3) & 41B aktywnych (Large 3); Apache 2.0 & Model europejski, w pełni otwarta licencja Apache 2.0; dobry punkt odniesienia dla modeli ,,frontier'' przy zachowaniu reprodukowalności. \\ \hline
DeepSeek V4 & MoE & Silny model do długiego kontekstu; reprezentuje klasę dużych modeli MoE dostępnych do uruchomienia lub przez API uczelni. \\ \hline
\end{tabularx}

\subsection*{Oś lokalna -- modele uruchamialne na jednej karcie GPU}

Modele tej osi służą do zbadania tezy, że rekomendacja LLM jest wykonalna bez infrastruktury chmurowej -- co ma bezpośrednie znaczenie praktyczne dla mniejszego sklepu e-commerce.

\begin{tabularx}{\linewidth}{|>{\bfseries\color{primary}}p{3.0cm}|p{2.2cm}|X|}
\hline
\rowcolor{primary}
\textcolor{white}{\textbf{Model}} & \textcolor{white}{\textbf{Skala}} & \textcolor{white}{\textbf{Uzasadnienie doboru}} \\ \hline
Mistral Small & 24B & Mieści się w 16 GB VRAM po kwantyzacji 4-bit; dobry kompromis jakość/rozmiar dla wariantu lokalnego. \\ \hline
Gemma 4 & ok. 27B & Rodzina Google o silnych wynikach na zadaniach językowych; alternatywa dla Mistral Small w porównaniu modeli lokalnych. \\ \hline
Phi-4 & 14B & Mały, ,,gęsty'' model o dużej efektywności -- reprezentuje dolną granicę skali; mieści się w VRAM po kwantyzacji 8-bit. \\ \hline
\end{tabularx}

\subsection*{Oś polska -- punkt odniesienia i lokalny wkład pracy}

\textbf{Bielik} to seria otwarto-źródłowych modeli językowych dla języka polskiego, rozwijana przez społeczność SpeakLeash przy wsparciu Akademii Górniczo-Hutniczej i innych uczelni. Modele Bielik v3 (7B i 11B) powstały przez kontynuację treningu (\emph{continual pre-training}) na bazie modelu Llama, z użyciem korpusu polskich tekstów liczącego dziesiątki miliardów tokenów. Innowacją jest autorski tokenizer zoptymalizowany pod język polski, redukujący liczbę tokenów potrzebnych do reprezentacji polskiego tekstu -- co przekłada się na krótszy czas inferencji. Modele są dostępne na licencji Apache 2.0.

\textbf{PLLuM} (Polish Large Language Model) to inicjatywa finansowana ze środków publicznych, realizowana przez konsorcjum obejmujące OPI PIB, Politechnikę Gdańską i inne jednostki. PLLuM oferuje modele w rozmiarach od 8B do 70B parametrów, trenowane lub dostrajane na polskojęzycznych danych ogólnych i specjalistycznych.

W temacie D2 -- zgodnie z ustaleniem językowym -- modele polskie pełnią rolę \emph{punktu odniesienia}: pozwalają sprawdzić, czy model trenowany na polskich danych daje przewagę przy polskojęzycznych opisach produktów, mimo że prompty operacyjne sformułowane są po angielsku. Hipoteza H3 zakłada, że Bielik v3 11B dorównuje wielokrotnie większym modelom uniwersalnym przy ułamku kosztu -- jest to teza interesująca naukowo i o wyraźnym lokalnym charakterze.

\begin{infobox}{Uwaga praktyczna -- dobór końcowy}
Do badania nie trzeba używać wszystkich wymienionych modeli. Rekomendowany zestaw roboczy to \textbf{4--6 modeli}: dwa z osi głównej (np. Qwen 3.5 i Mistral), jeden--dwa lokalne (np. Mistral Small lub Phi-4) oraz co najmniej jeden polski (Bielik v3 11B). Taki dobór pokrywa pełne spektrum skali i pochodzenia, a jednocześnie utrzymuje liczbę warunków eksperymentu na poziomie wykonalnym -- co jest istotne dla mocy statystycznej testów (patrz rozdział~\ref{sec:liczebnosc}).
\end{infobox}

\newpage

% =============================================================================
% 6. TEMAT P4
% =============================================================================
\section{Temat P4 -- pełny opis}

\topicheader{Temat P4}{promotor}{Wnioskowanie przedziałowe (IT2-FLS) a klasyczny system rozmyty -- z opcjonalnym rozszerzeniem o hybrydę neuronowo-rozmytą ANFIS}

\topicsection{Cel i pytanie badawcze}
Temat stanowi rozszerzenie istniejącego w aplikacji modułu wnioskowania rozmytego typu pierwszego (system hybrydowy Mamdani-TSK, klasa \texttt{SimpleFuzzyInference} -- doprecyzowanie klasyfikacji poniżej). Pytanie badawcze: czy formalne modelowanie niepewności (logika rozmyta typu drugiego) realnie poprawia jakość rekomendacji w warunkach rzadkich danych, czy jedynie zwiększa koszt obliczeniowy bez wymiernych korzyści. \textbf{Trzon pracy} stanowi porównanie \emph{dwóch} systemów: typu pierwszego (baza odniesienia) oraz typu drugiego przedziałowego (IT2-FLS). Jako \textbf{rozszerzenie warunkowe} -- jeśli pozwoli na to czas i dostępność danych etykietowanych -- badanie obejmuje trzeci system: hybrydę neuronowo-rozmytą ANFIS (uzasadnienie i ograniczenia opisano w osobnej sekcji poniżej).

\topicsection{Punkt wyjścia -- istniejący system rozmyty}
W obecnej implementacji (\texttt{fuzzy\_logic\_engine.py}) system operuje na trzech zmiennych językowych: \emph{cena} (zbiory: tani / średni / drogi -- funkcje trójkątne i trapezowe), \emph{jakość} (niska / średnia / wysoka, na podstawie oceny 1--5) oraz \emph{popularność} (niska / średnia / wysoka, na podstawie liczby wyświetleń). Baza reguł liczy sześć reguł (np. ,,R1: wysoka jakość w dobrej cenie''). Faza wnioskowania jest zgodna z podejściem Mamdaniego -- rozmyte poprzedniki (przesłanki) oraz operatory T-normy (minimum) i S-normy (maksimum). Etap defuzyfikacji nie definiuje jednak rozmytych zbiorów dla zmiennej wyjściowej: wynik liczony jest jako średnia ważona stałych wartości liczbowych w następnikach reguł.

\begin{infobox}{Doprecyzowanie terminologiczne -- Mamdani czy TSK}
Ponieważ następniki reguł są stałymi wartościami liczbowymi (a nie zbiorami rozmytymi), obecny system jest formalnie \textbf{systemem wnioskowania typu Sugeno zerowego rzędu (TSK-0)}, a nie klasycznym systemem Mamdaniego. Zgodnie z Ross (2010, \textit{Fuzzy Logic with Engineering Applications}) model Sugeno rzędu zerowego, w którym następnik każdej reguły jest stałą, bywa traktowany jako szczególny przypadek systemu Mamdaniego z następnikami w postaci singletonów rozmytych. Z uwagi na to strukturalne podobieństwo (rozmyte poprzedniki, struktura IF-THEN, operatory T-/S-norm w stylu Mamdaniego) najprecyzyjniejszym określeniem jest \textbf{system hybrydowy Mamdani-TSK} (lub: ,,system wnioskowania typu Sugeno zerowego rzędu''). Nieścisłość dotyczy wyłącznie klasyfikacji teoretycznej -- sama implementacja działa poprawnie zgodnie z teorią wnioskowania rozmytego.
\end{infobox}

Niezależnie od klasyfikacji Mamdani/TSK, obecny system jest systemem \emph{typu pierwszego} (stopnie przynależności to ostre liczby) -- i to właśnie ta cecha (a nie typ następnika) jest osią rozszerzenia w niniejszym temacie. System ten jest bezpośrednim punktem odniesienia dla nowych wariantów -- nie wymaga przebudowy, służy jako podstawa (baseline).

\topicsection{Badane systemy -- trzon i rozszerzenie}
\textbf{Trzon pracy (dwa systemy):}
\begin{enumerate}
\item \textbf{System typu pierwszego (T1, Mamdani-TSK)} -- istniejący system opisany wyżej. Stopień przynależności jest pojedynczą, ostrą liczbą -- np. cena pewnego produktu należy do zbioru ,,tani'' w stopniu dokładnie 0,5.

\item \textbf{System typu drugiego przedziałowy (IT2-FLS)} -- modeluje niepewność samej funkcji przynależności. Zamiast jednej liczby stopień przynależności jest \emph{przedziałem} -- np. od 0,4 do 0,6. Każdy zbiór rozmyty typu drugiego charakteryzowany jest parą funkcji wyznaczających obszar niepewności (\emph{Footprint of Uncertainty}, FOU). Lepiej oddaje to rzeczywistą niepewność danych e-commerce: ceny zmieniają się dynamicznie, oceny są subiektywne, opinii bywa mało. Implementacja korzysta z biblioteki \texttt{pyit2fls}, dostarczającej klasy zbiorów IT2FS, systemy Mamdani i TSK oraz algorytmy redukcji typu i defuzyfikacji Karnika-Mendela.
\end{enumerate}

Para T1 kontra IT2-FLS jest osią badania: oba systemy korzystają z tej samej, eksperckiej bazy reguł, różnią się wyłącznie sposobem reprezentacji niepewności -- co czyni porównanie czystym i jednoznacznym (żaden z systemów nie wymaga danych etykietowanych do treningu).

\textbf{Rozszerzenie warunkowe (trzeci system):}
\begin{enumerate}\setcounter{enumi}{2}
\item \textbf{Hybryda neuronowo-rozmyta ANFIS} -- system, w którym parametry funkcji przynależności i wagi reguł nie są ustalane ręcznie przez eksperta, lecz \emph{uczone} z danych algorytmem uczenia hybrydowego. ANFIS pozwoliłby porównać trzy paradygmaty wiedzy: stałą wiedzę ekspercką (T1), wiedzę ekspercką z modelowaniem niepewności (IT2) oraz wiedzę uczoną z danych (ANFIS). ANFIS wnosi jednak istotne ograniczenia metodologiczne (problem etykiet treningowych, ryzyko uczenia się danych syntetycznych), dlatego traktowany jest jako rozszerzenie warunkowe, a nie równoprawny trzon -- szczegóły w sekcji ,,Analiza wykonalności ANFIS''.
\end{enumerate}

\topicsection{Zakres prac implementacyjnych}
Powstaje nowy moduł wnioskowania typu drugiego, działający równolegle do istniejącego modułu typu pierwszego; moduł ANFIS jest realizowany warunkowo.

\begin{itemize}
\item \textbf{Moduł IT2-FLS} (trzon) -- definicja zbiorów rozmytych typu drugiego dla wszystkich pojęć językowych systemu (cena, jakość, popularność), każdy z wyznaczonym obszarem niepewności FOU; przepisanie istniejącej bazy sześciu reguł do postaci typu drugiego. Ponieważ system bazowy jest typu Sugeno zerowego rzędu (Mamdani-TSK), naturalnym odpowiednikiem typu drugiego jest wariant IT2-TSK (klasa \texttt{IT2TSK} z \texttt{pyit2fls}); dla porównania rozważany jest również wariant IT2-Mamdani (klasa \texttt{IT2Mamdani}). Redukcja typu i defuzyfikacja -- algorytmem Karnika-Mendela.
\item \textbf{Panel diagnostyczny} (trzon) -- rozszerzenie istniejącego panelu logiki rozmytej w interfejsie administratora o wizualizację obszaru niepewności FOU, aktywacji reguł górnych i dolnych oraz przedziału wyjściowego.
\item \textbf{Mechanizm porównania równoległego} (trzon) -- uruchomienie porównywanych systemów na tych samych danych wejściowych i zapis wyników do wspólnej tabeli eksperymentów.
\item \textbf{Moduł ANFIS} (rozszerzenie warunkowe) -- struktura sieci adaptacyjnej odwzorowująca system rozmyty, z parametrami funkcji przynależności uczonymi z danych. Realizowany pod warunkiem rozstrzygnięcia kwestii etykiet treningowych opisanej w sekcji ,,Analiza wykonalności ANFIS''.
\end{itemize}

\topicsection{Stack technologiczny}
\begin{stacktab}
\stackrow{System typu drugiego}{biblioteka \texttt{pyit2fls}}
\stackrow{ANFIS}{implementacja na bazie PyTorch / TensorFlow lub dedykowanej biblioteki ANFIS}
\stackrow{Numeryka}{NumPy, SciPy}
\stackrow{Wizualizacja FOU}{matplotlib -- wykresy obszaru niepewności}
\stackrow{Backend}{Django -- nowe moduły logiki rozmytej obok istniejącego \texttt{fuzzy\_logic\_engine.py}}
\stackrow{Walidacja przedziału wyjściowego}{symulacja Monte Carlo -- perturbacja wejść w zakresie ich niepewności i sprawdzenie, czy rozrzut wyników mieści się w przedziale wyznaczonym przez system typu drugiego (pokrycie / coverage)}
\end{stacktab}

\topicsection{Innowacyjność}
Wkład pracy obejmuje trzy aspekty. Po pierwsze, jedno z nielicznych zastosowań logiki rozmytej typu drugiego w kontekście systemów rekomendacji w polskim środowisku akademickim -- większość prac dotyczących IT2-FLS dotyczy sterowania, robotyki i finansów. Po drugie, empiryczna weryfikacja pytania, czy modelowanie niepewności faktycznie poprawia jakość rekomendacji w warunkach rzadkich danych. Po trzecie, propozycja nowej metryki -- szerokości przedziału wyjściowego (niepewności rozmytej) rekomendacji -- jako miary pewności systemu; otwiera to drogę do rankingu świadomego pewności. Opcjonalne rozszerzenie o ANFIS dokłada czwarty wątek: porównanie wiedzy eksperckiej z wiedzą uczoną z danych.

\topicsection{Analiza wykonalności ANFIS -- dlaczego rozszerzenie warunkowe}
Zasugerowane włączenie ANFIS jako jednej z możliwości. Po analizie obecnej aplikacji ANFIS jest \emph{wykonalny}, lecz wnosi tarcia, których para T1 kontra IT2-FLS nie ma -- stąd jego rola jako rozszerzenia warunkowego, a nie równoprawnego trzonu.

\textbf{Co przemawia za ANFIS.} Pod względem architektury ANFIS pasuje do tego systemu \emph{dobrze}: skoro obecny silnik jest typu Sugeno zerowego rzędu (Mamdani-TSK), to ANFIS -- który natywnie używa następników w stylu Sugeno -- jest jego bezpośrednim, uczonym odpowiednikiem. Model jest mały (3 wejścia, kilkadziesiąt parametrów), trenuje się offline w sekundy, nie wymaga GPU, a wytrenowane wagi można materializować w bazie tak jak inne tabele pochodne. Od strony platformy e-commerce (Django/React) nie ma żadnych przeszkód.

\textbf{Główne ograniczenie -- problem etykiet treningowych.} ANFIS to model uczony nadzorowany: potrzebuje par ,,wejście $\rightarrow$ poprawne wyjście''. Wejścia są dostępne (cena, jakość, popularność, dopasowanie kategorii), ale obecny system rozmyty \emph{nie definiuje} docelowej wartości wyjściowej -- score jest wyznaczany regułami eksperckimi, nie ma ,,prawdy'', do której ANFIS miałby się dopasować. Trzeba więc świadomie wybrać cel uczenia, a każdy wariant ma wadę:
\begin{itemize}
\item \emph{Cel = przyszły zakup lub wysoka ocena} -- sensowny i obronny, ale wtedy ANFIS staje się klasyfikatorem dopasowania konkurującym raczej z naiwnym Bayesem czy filtracją kolaboratywną niż z logiką rozmytą; pojawia się pytanie ,,czemu ANFIS, a nie zwykły gradient boosting''.
\item \emph{Cel = score obecnego systemu rozmytego} -- ANFIS uczy się naśladować T1, więc z definicji go nie przewyższy (T1 jest jego nauczycielem); naukowo słabe.
\item \emph{Cel = oceny opinii (rating)} -- najbliższe duchowi systemu, ale opinie są rzadkie, a to właśnie scenariusz, w którym ANFIS radzi sobie najgorzej.
\end{itemize}

\textbf{Drugie ograniczenie -- ryzyko uczenia się danych syntetycznych.} Jeśli ANFIS trenowany jest na badawczym seederze, który sam zawiera wbudowaną logikę preferencji, istnieje ryzyko, że model uczy się seederа, a nie zjawiska. Para T1 kontra IT2-FLS tego problemu nie ma w ogóle, bo nie jest uczona z danych.

\begin{infobox}{Decyzja projektowa -- ANFIS jako plan B}
Trzon tematu P4 to czyste, jednoznaczne porównanie \textbf{T1 (Mamdani-TSK) kontra IT2-FLS} -- oba na tej samej bazie reguł eksperckich, bez potrzeby etykiet, bez ryzyka uczenia się seederа. ANFIS zostaje oznaczony jako \textbf{rozszerzenie warunkowe}
\end{infobox}

\topicsection{Hipotezy badawcze}
\begin{itemize}
\item \textbf{H1.} System typu drugiego osiąga wyższą metrykę Precision@K niż system typu pierwszego dla użytkowników z małą liczbą ocen -- w warunkach rzadkich danych modelowanie niepewności ma większe znaczenie.
\item \textbf{H2.} Szerokość przedziału wyjściowego obliczonego przez system typu drugiego koreluje istotnie z błędem predykcji -- szeroki przedział wskazuje rzeczywistą niepewność systemu.
\item \textbf{H3.} Koszt obliczeniowy systemu typu drugiego (zwłaszcza defuzyfikacji Karnika-Mendela) jest istotnie wyższy niż systemu typu pierwszego, lecz opłacalny tylko w wybranych scenariuszach (rzadkie dane, nowi użytkownicy).
\item \textbf{H4} (warunkowa, dotyczy rozszerzenia ANFIS). ANFIS, ucząc parametry z danych, dorównuje lub przewyższa system typu pierwszego o stałych parametrach przy celu uczenia ustawionym na przyszły zakup, a jego jakość rośnie wraz z liczbą próbek treningowych, osiągając plateau.
\end{itemize}

\topicsection{Plan eksperymentu}
\textbf{Trzon.} Wykorzystanie istniejącego modułu typu pierwszego jako bazy odniesienia. Implementacja wariantu IT2-FLS z obszarem niepewności wyznaczonym jako odchylenie procentowe od funkcji typu pierwszego (parametr regulowany). Test na trzech rozłącznych podzbiorach użytkowników: o niskiej, średniej i wysokiej liczbie ocen. Porównanie systemu typu pierwszego, drugiego oraz wariantu hybrydowego (IT2 dla użytkowników o ubogiej historii, T1 dla pozostałych). Walidacja przedziału wyjściowego metodą Monte Carlo -- perturbacja wejść w zakresie ich niepewności i sprawdzenie pokrycia (coverage) przedziału wyznaczonego przez system typu drugiego. Pomiar opóźnienia obliczeń dla każdego z systemów. Analiza statystyczna -- test Friedmana z testem post-hoc (gdy porównywane są więcej niż dwa systemy) lub test Wilcoxona dla par powiązanych (dla pary T1 kontra IT2).

\textbf{Rozszerzenie warunkowe.} Jeśli realizowany będzie moduł ANFIS: implementacja i trening z celem uczenia ustawionym na przyszły zakup, z jawną kontrolą poziomu szumu seederа oraz walidacją krzyżową na niezależnym zbiorze. ANFIS dołącza wówczas jako trzeci system do porównania opisanego wyżej.

\begin{literatura}
\item Mendel, J. M. (2017). \textit{Uncertain Rule-Based Fuzzy Systems: Introduction and New Directions}, wyd. 2. Springer.
\item Karnik, N. N., Mendel, J. M. (2001). Centroid of a type-2 fuzzy set. \textit{Information Sciences}, 132(1-4), 195--220.
\item Wu, D. (2013). Approaches for reducing the computational cost of interval type-2 fuzzy logic systems. \textit{IEEE Transactions on Fuzzy Systems}, 21(1), 80--99.
\item Haghrah, A. A., Ghaemi, S. (2023). PyIT2FLS: A New Python Toolkit for Interval Type 2 Fuzzy Logic Systems. \textit{arXiv:1909.10051}.
\item Jang, J.-S. R. (1993). ANFIS: Adaptive-Network-Based Fuzzy Inference System. \textit{IEEE Transactions on Systems, Man, and Cybernetics}, 23(3), 665--685.
\item Chen, C., John, R. i in. Type-1 and Interval Type-2 ANFIS: A Comparison. \textit{Materiały konferencyjne IEEE}.
\item Castillo, O., Melin, P. (2008). \textit{Type-2 Fuzzy Logic: Theory and Applications}. Springer.
\end{literatura}

\newpage

% =============================================================================
% 7. LICZEBNOŚĆ PRÓBY
% =============================================================================
\section{Wiarygodność statystyczna i forma ewaluacji}
\label{sec:liczebnosc}

Wiarygodność (istotność) statystyczna testów jest dyktowana minimalnymi wymaganiami stosowanych testów oraz użytkowników. Poniżej przegląd literatury potwierdzający i precyzujący.

\begin{infobox}{Forma ewaluacji w obu pracach}
Główną i rozstrzygającą formą oceny w obu pracach jest \textbf{ewaluacja offline} -- automatyczny pomiar metryk rankingowych (Precision@K, Recall@K, nDCG@K) na badawczym zbiorze danych, z testami statystycznymi po użytkownikach. \textbf{Nie planujemy badania z udziałem użytkowników}. Jako jakościowe uzupełnienie przewidujemy \textbf{ocenę ekspercką} przeprowadzaną przez nas - autorów na niewielkiej, kontrolnej próbie przypadków. Poniższy przegląd literatury służy zatem dwóm celom: (1) uzasadnia liczebność \emph{użytkowników w zbiorze danych} dla ewaluacji offline oraz (2) pokazuje, dlaczego pełnoskalowe badanie z udziałem ludzi nie jest dla tej pracy konieczne.
\end{infobox}

\subsection*{Liczebność a istotność statystyczna -- liczby z literatury}

\begin{itemize}
\item \textbf{Nielsen Norman Group} (przegląd zaleceń dotyczących liczebności prób w badaniach ilościowych) wskazuje, że \textbf{40 uczestników} jest liczbą odpowiednią dla większości badań ilościowych -- przy poziomie ufności 95\% i marginesie błędu 15\%. Liczbę można obniżyć do \textbf{30} przy poziomie ufności 90\% lub do \textbf{20--25} w projektach o ograniczonym budżecie, choć autorzy określają to jako podejście ryzykowne. Wniosek operacyjny tego samego źródła: \emph{aby uzyskać statystycznie istotne wyniki, należy przebadać co najmniej 20 użytkowników; ciaśniejsze przedziały ufności wymagają większej liczby}.
\\ \textit{Źródło: \url{https://www.nngroup.com/articles/summary-quant-sample-sizes/}}

\item \textbf{Knijnenburg \& Willemsen}, \emph{Evaluating Recommender Systems with User Experiments} -- praca metodyczna wprost poświęcona eksperymentom z użytkownikami w systemach rekomendacji. Autorzy zwracają uwagę, że jedno z pierwszych badań user-centric systemu rekomendacji objęło zaledwie \textbf{19 uczestników} i mimo wykrycia istotnych efektów było \emph{niedostatecznie zasilone} (\emph{underpowered}) -- wniosków nie dało się uogólnić. Kluczowa wskazówka metodyczna: liczba uczestników potrzebna dla danej mocy statystycznej rośnie \emph{liniowo z liczbą warunków eksperymentu}, dlatego należy utrzymywać liczbę porównywanych warunków na niskim poziomie.
\\ \textit{Źródło: \url{https://www.usabart.nl/portfolio/KnijnenburgWillemsen-UserExperiments.pdf}}

\item \textbf{Nielsen Norman Group}, \emph{How Many Test Users in a Usability Study} -- klasyczne rozróżnienie: badania \emph{jakościowe} (wykrywanie problemów) wymagają około 5 użytkowników, natomiast badania \emph{ilościowe} nastawione na statystykę wymagają co najmniej 20 użytkowników.
\\ \textit{Źródło: \url{https://www.nngroup.com/articles/how-many-test-users/}}

\item Przegląd metodyczny w czasopismach medycznych (\emph{Sample size: how many participants do I need in my research?}) podkreśla, że dokładna liczebność wynika z planowanej \emph{mocy statystycznej} (zwykle 80\%), zakładanej wielkości efektu i poziomu istotności ($\alpha = 0{,}05$) -- liczbę należy wyznaczyć analizą mocy, a nie przyjmować arbitralnie.
\\ \textit{Źródło: \url{https://pmc.ncbi.nlm.nih.gov/articles/PMC4148275/}}
\end{itemize}

\subsection*{Wnioski dla obu prac}

\begin{okbox}{Rekomendacja dla obu prac}
\textbf{Ewaluacja offline -- forma główna (metryki rankingowe, testy statystyczne po użytkownikach).} Badawczy zbiór danych powinien obejmować rzędu \textbf{150--300 użytkowników} z odpowiednio bogatą historią -- jest to liczba swobodnie osiągalna przez parametryzowany seeder, a jednocześnie wystarczająca, by testy po użytkownikach (Friedman, Wilcoxon) miały wysoką moc. Tu liczebność nie jest ograniczeniem -- ogranicza ją wyłącznie realizm danych, nie ich pozyskanie. To na tej ewaluacji opiera się główny ciężar dowodowy obu prac.

\textbf{Ocena ekspercka -- uzupełnienie jakościowe.} W miejsce niemożliwego do przeprowadzenia badania z udziałem użytkowników stosujemy ocenę ekspercką: ocenimy jakość rekomendacji na \textbf{niewielkiej, kontrolnej próbie przypadków} (np. kilkadziesiąt zapytań rekomendacyjnych obejmujących różne persony i scenariusze, w tym zimny start). Ocena ta nie służy do wyciągania wniosków statystycznych o populacji użytkowników -- pełni rolę \emph{walidacji jakościowej}: weryfikuje, czy rekomendacje wysoko ocenione przez metryki są też sensowne ,,po ludzku'', oraz wychwytuje błędy, których metryki nie pokazują (np. rekomendacje trafne, lecz nietrafione kontekstowo, albo -- w temacie D2 -- halucynacje modelu). Aby ograniczyć subiektywizm, ocena prowadzona jest według ustalonej rubryki (np. trafność, różnorodność, poprawność opisu) i, gdzie to możliwe, w trybie ślepym względem metody.

\textbf{Liczba porównywanych warunków.} Niezależnie od formy oceny liczbę porównywanych systemów utrzymujemy nisko -- rekomendowany zestaw 4--6 modeli w temacie D2 oraz 2 systemy w trzonie tematu P4 (T1 kontra IT2-FLS), ewentualnie 3 z opcjonalnym ANFIS.

\textbf{Zasada nadrzędna.} Wiarygodność wyników opiera się na ewaluacji offline; ostateczne parametry testów statystycznych (Friedman, Wilcoxon) zostaną dobrane analizą mocy przy założonej mocy 80\% i $\alpha = 0{,}05$ i opisane w rozdziale metodycznym każdej z prac. Ocena ekspercka stanowi jakościowe potwierdzenie, nie odrębne badanie statystyczne.
\end{okbox}

% =============================================================================
% 8. PODZIAŁ PRAC
% =============================================================================
\section{Podział prac i harmonogram ramowy}

\begin{tabularx}{\linewidth}{|>{\bfseries\color{primary}}p{3.4cm}|X|}
\hline
\rowcolor{primary}
\textcolor{white}{\textbf{Zakres}} & \textcolor{white}{\textbf{Realizacja}} \\ \hline
Wspólnie & Badawczy seeder \texttt{seed\_research}, polecenie \texttt{prepare\_research\_db}, wspólny moduł metryk, model \texttt{ExperimentRun}, wspólny rozdział zestawiający wyniki obu prac. \\ \hline
Dawid Olko (D2) & Moduł rekomendacji LLM (prompting, embeddings, RAG), dobór i konfiguracja modeli, eksperymenty porównawcze LLM kontra metody klasyczne, ocena halucynacji i kosztu, badanie modeli polskich. \\ \hline
Piotr Smoła (P4) & Moduł IT2-FLS (trzon), przepisanie bazy reguł rozmytych, panel diagnostyczny FOU, eksperymenty porównawcze T1 kontra IT2-FLS, walidacja przedziałów Monte Carlo; opcjonalnie moduł ANFIS jako rozszerzenie warunkowe. \\ \hline
\end{tabularx}

Każda praca ma własny tytuł i własne rozdziały. Wstęp obu prac może zawierać wspólny akapit o tym samym repozytorium i tym samym badawczym zbiorze danych. Ramowy harmonogram: (1) fundament wspólny; (2) implementacja modułów własnych; (3) eksperymenty i zbieranie wyników; (4) analiza statystyczna i redakcja; (5) wspólny rozdział zestawiający wyniki.

\vspace{1em}

\begin{infobox}{Podsumowanie}
Tematy D2 i P4 zostały wstępnie zatwierdzone. Niniejszy dokument uszczegóławia ich zakres w oparciu o faktyczną analizę aplikacji \emph{SmartRecommender} oraz uwzględnia wszystkie cztery wskazówki: język pracy (polski) i promptów (angielski), polskie modele Bielik i PLLuM jako punkt odniesienia w temacie D2, możliwość rozszerzenia tematu P4 o ANFIS (po analizie wykonalności potraktowany jako rozszerzenie warunkowe obok trzonu T1 kontra IT2-FLS) oraz oparcie liczebności prób na przeglądzie literatury. Oba tematy współdzielą fundament badawczy, co umożliwia bezpośrednie zestawienie wyników i -- w razie interesujących rezultatów -- wspólną publikację.
\end{infobox}

\end{document}
